% Main Figures and Tables

% Table 1: Performance Metrics
\begin{table}[htbp]
  \centering
  \caption{\textbf{Tissue-specific classification performance.} Balanced accuracy (BA) is reported with 95\% bootstrap confidence intervals. MAE is reported in stage units. A dash indicates that ordinal MAE is not defined for two-class models.}
  \label{tab:performance}
  \sisetup{round-mode = places, round-precision = 3}
  \begin{tabular}{lllll}
    \toprule
    Tissue & Scheme & BA [95\% CI] & Macro F1 & MAE \\
    \midrule
    Muscle & 4-class & 0.887~[0.833, 0.934] & 0.888 & 0.095 \\
    Brain & 4-class & 0.905~[0.846, 0.952] & 0.894 & 0.157 \\
    Liver & 4-class & 0.924~[0.870, 0.969] & 0.919 & 0.081 \\
    Blood & 3-class & 0.821~[0.699, 0.934] & 0.762 & 0.163 \\
    Small intestine & 3-class & 0.858~[0.763, 0.939] & 0.806 & 0.291 \\
    Lung & 3-class & 0.867~[0.746, 0.970] & 0.874 & 0.111 \\
    Adipose & 2-class & 0.907~[0.808, 0.982] & 0.907 & -- \\
    Testis & 2-class & 0.955~[0.850, 1.000] & 0.952 & -- \\
    \bottomrule
  \end{tabular}
\end{table}

% Figure 1: Machine Learning Pipeline (Previously Supplementary Figure 1)
\begin{figure}[htbp]
  \centering
  \includegraphics[width=\textwidth]{supplementary_figure1_ml_pipeline.pdf}
  \caption{\textbf{Machine learning pipeline schematic.} End-to-end workflow from raw PigGTEx transcriptomes through preprocessing, adaptive stage granularity assignment, model training, calibration, and evaluation outputs.}
  \label{fig:mlpipeline}
\end{figure}

% Figure 2: Study Overview (Previously Figure 1)
\begin{figure}[htbp]
  \centering
  \includegraphics[width=\textwidth]{figure1_methodology.pdf}
  \caption{\textbf{Overview of the developmental staging workflow.} (a) Age bins used to harmonise PigGTEx metadata into five developmental stages. (b) Tissue-level sample availability and resulting stage granularity assignments. Bottom panel: End-to-end machine learning pipeline from data ingestion through calibrated evaluation.}
  \label{fig:methodology}
\end{figure}

% Figure 3: Performance (Previously Figure 2)
\begin{figure}[htbp]
  \centering
  \includegraphics[width=\textwidth]{figure2_performance.pdf}
  \caption{\textbf{Classifier performance diagnostics.} (a) Balanced accuracy, macro F1, mean absolute error, and Spearman correlations with 95\% bootstrap intervals. (b) Confusion matrices emphasising near-miss errors. (c) Correlation between predicted stage scores and chronological age across representative tissues.}
  \label{fig:performance}
\end{figure}

% Figure 4: Molecular Signatures (Previously Figure 3)
\begin{figure}[htbp]
  \centering
  \includegraphics[width=\textwidth]{figure3_molecular_signatures.pdf}
  \caption{\textbf{Molecular signatures of porcine developmental staging.} (a) Tissue-specific UMAP embeddings coloured by developmental stage. (b) Feature overlap among the top 50 markers per tissue. (c) Exemplar marker trajectories across developmental stages for key tissue-specific genes.}
  \label{fig:signatures}
\end{figure}

% Figure 5: Biological Validation (Previously Figure 4)
\begin{figure}[htbp]
  \centering
  \includegraphics[width=\textwidth]{figure4_biological_validation.pdf}
  \caption{\textbf{Functional validation of stage-informative genes.} (a) Tissue-resolved gene ontology enrichments for top biological process terms. (b) Network view of developmental pathway modules. (c) Expression tracking of canonical markers across stages, illustrating biological concordance with classifier predictions.}
  \label{fig:biovalidation}
\end{figure}